\subsection{Relation to MOND-like Phenomenology}
  \label{subsec:mond_relation}

  The phenomenology discussed in this work shares qualitative similarities with
  acceleration-based frameworks such as Modified Newtonian Dynamics (MOND),
  originally introduced to account for flat galaxy rotation curves without
  invoking dark matter~\cite{Milgrom1983,Famaey2012Review}.

  In particular, both approaches identify a characteristic acceleration scale
  below which Newtonian expectations fail.
  However, the conceptual status of this scale differs fundamentally.

  Three key observational discriminants distinguish the present framework from MOND:
  \begin{itemize}
    \item \textbf{Environment-dependent Hubble measurements}: The effective framework predicts a correlation between
    locally inferred \(H_0\) and large-scale environment (Eq.~\eqref{eq:hloc}), absent in MOND where \(a_0\)
    is universal and does not couple to cosmological expansion.

    \item \textbf{Redshift dependence of the Hubble offset}: The deviation \(\Delta H/H\)
    is predicted to converge to zero at high redshift (Section~\ref{subsec:redshift_dependence}
    ), while MOND provides no mechanism for redshift-dependent modifications of \(H_0\).

    \item \textbf{Ultra-diffuse galaxy kinematics}: UDGs in voids are predicted to exhibit systematically lower
    asymptotic velocities than UDGs in clusters(Section~\ref{subsec:udg}), whereas MOND predicts a universal
    \(v_{\infty} \propto (M a_0)^{1/4}\) independent of environment.
  \end{itemize}
  These differences arise because the saturation scale \(a_\star\) in the effective framework is tied to the
  \textbf{projection efficiency}
  of the underlying geometric substrate (Section 4.3 of~\cite{Beau2026c}), not to a modified force law.

  In the present framework, the low-acceleration behavior is not postulated as a modification of gravity or inertia.
  It arises as a saturation of the effective transition resolution in low-density environments.
  The scale $a_\star$ characterizes a limit of descriptive refinement, not a new dynamical constant.

  This distinction becomes especially important in cosmological applications.
  MOND-based extensions encounter well-known difficulties when applied consistently to large-scale structure
  and expansion observables.
  By contrast, the effective saturation framework preserves the standard background expansion
  and modifies only the local inference of kinematic quantities in diffuse late-time environments.
